\documentclass[11pt]{article}
\usepackage[margin=1in]{geometry}
\usepackage{xcolor}
\usepackage{booktabs}
\usepackage{enumitem}
\usepackage[hidelinks]{hyperref}
\usepackage{times}

\newcommand{\rcomment}[1]{\begin{quote}\itshape\color{black!70}#1\end{quote}}
\newcommand{\resp}{\noindent\textbf{Response:} }
\newcommand{\excerpt}[1]{\begin{quote}\small\color{blue!70!black}#1\end{quote}}
\setlength{\parskip}{4pt}

\title{Point-by-Point Response to the Editor and Reviewers\\
\large Revision of Submission b3507973-de3c-4a56-a992-03d07bc762ee\\
``Scaling vs.\ Architecture: An Evaluation of SLMs for Automated Docstring Generation''\\
(revised title: ``Capacity vs.\ Architecture: An Evaluation of SLMs for Automated Docstring Generation'')}
\author{Balaji Venktesh, Amsaprabhaa M, Gireesh Sundaram}
\date{}

\begin{document}
\maketitle

\noindent Dear Editor and Reviewers,

We would like to express our sincere gratitude for the time and effort you dedicated to reviewing our manuscript. The comments were constructive, precise, and valuable, and they have materially improved both the rigor and the honesty of this work. We have studied every comment carefully and made detailed revisions, summarized below and answered point by point. For easy cross-referencing, all changed or added material in the marked-up manuscript copy is set in blue font.

\textit{Notation.} Comments are tagged R1.$n$/R2.$n$/R3.$n$ for the $n$-th comment of each reviewer, in the order raised. Where the revised manuscript is quoted, excerpts appear in indented blue text so changes can be verified without flipping back to the paper.

\medskip
\noindent Best regards,\\
The Authors

\section*{To the Editor}

We thank the editor and all three reviewers for their careful and demanding reviews. The instruction to ``ensure the results are accurately reported'' prompted us to go beyond the specific comments: we re-verified every experimental claim in the manuscript against our execution logs, code, and released artifacts. This audit confirmed the reviewers' concerns and surfaced additional inaccuracies, all of which are corrected in this revision. In summary:

\begin{itemize}[leftmargin=1.4em]
  \item \textbf{Corrected reporting.} The evaluated models are Llama~3.2~\textbf{3B} (Ollama tag \texttt{llama3.2:latest}; the submitted manuscript incorrectly said ``8B''), Qwen3~8B, and Phi-4~14B; the abstract's ``Qwen2.5'' was a pilot-phase remnant. The benchmark comprises \textbf{67 curated classes} (2{,}613 core generations), not 1{,}400/54{,}600 as stated; the erroneous figures arose from conflating a planned corpus size with the executed benchmark, and every count, table, and library-distribution figure has been corrected to the executed data. Exact model tags, dataset composition, and hardware are now stated precisely (revised \S4).
  \item \textbf{Re-analysis.} All faithfulness scores were recomputed against a single reference (source code) for all strategies (R2.2), and all inferential statistics were re-derived from linear mixed-effects models at the per-class unit of analysis (R2.3).
  \item \textbf{New human validation.} The originally reported human--judge agreement ($r{=}0.925$) did not survive a strictly blind re-annotation and has been withdrawn; a blind two-annotator study replaces it and now anchors the paper's metric hierarchy (\S\ref{resp:r26}).
  \item \textbf{Six additional experiments} (summarized in \S\ref{sec:newexp}) address the reviewers' requests for a stronger few-shot baseline, a fairer retrieval corpus, a quantization control, gate-heuristic robustness, repeated-run variance, and a second LLM judge.
  \item \textbf{Rescoped claims.} The title adopts Reviewer~2's proposed wording; overstated conclusions (including our own headline finding on Few-Shot prompting, which the new ablation overturns) have been rewritten to match the evidence.
\end{itemize}

A marked-up copy of the revised manuscript with all changed or added material in blue is provided in the related-files section.

\section*{Summary of Changes: Consolidated Issue Map}

\begin{center}
\small
\begin{tabular}{p{0.4cm}p{3.6cm}p{2.4cm}p{7.6cm}}
\toprule
\# & \textbf{Issue} & \textbf{Sources} & \textbf{What we did} \\
\midrule
A & Model identities and reporting accuracy & R1.1, R2.1, Editor & Execution-log audit; all models corrected (Llama~3.2~\textbf{3B}, tag \texttt{llama3.2:latest}); benchmark scale corrected to the executed 67 classes / 2{,}613 generations; exact tags, composition, and hardware stated (\S4). Capacity finding re-verified and strengthened. \\
B & Commensurable faithfulness + correct unit of analysis & R2.2, R2.3 & All 2{,}613 outputs re-judged against source code ($k{=}3$ draws); all inference re-derived from mixed-effects models with class as random effect ($N{=}2{,}613$; ICC up to 0.28). New Table~5 replaces the Kruskal--Wallis/Mann--Whitney apparatus. \\
C & Judge validity and human validation & R1.4, R2.4, R2.6, R2 minors & Strictly blind two-annotator re-annotation (ICC $=0.63$); original $r{=}0.925$ withdrawn as anchoring artifact; second cross-family judge added; both judges fail blind validation ($r\le0.05$) while agreeing with each other ($\rho=0.979$); metric hierarchy rebuilt on validated metrics (new \S5.1). \\
D & Few-shot baseline strength & R2.8 & Three-arm exemplar ablation on all three models; dynamic matched exemplars beat zero-shot everywhere ($+0.04$ to $+0.09$); ``Few-Shot is harmful'' withdrawn; prompt defect disclosed (\S3.3, \S5.5). \\
E & Retrieval corpus fairness & R3.1 & Library-specific corpus ablation, all three models; corpus relevance matters directionally but gains stay $\le +0.012$ over zero-shot (\S5.3). \\
F & Quantization confound & R2.5 & FP16 vs.\ Q4 ablation; deltas $\le 0.01$, an order below effects of interest (\S5.4). \\
G & Critique-gate robustness & R1.2 & Gate description corrected; robust parser implemented (unit-tested); rerun on all three models; no conclusion depends on the gate (\S3.3, \S5.7, App.~H). \\
H & Scoping: title, reasoning labels, capacity claims & R2.7, R2.9, R3.3, R3.4 & Title adopts R2's wording; modes renamed ``reasoning-mode prompt variants'' with scope statement; capacity/training-quality framing with explicit design limitation; contributions sharpened (\S1, \S2.6, \S3.4, \S7.2). \\
I & Reasoning Tax mechanisms & R3.2 & New mechanism section: hallucination-by-inference, capacity-dependent format failures (all 5 runaway generations on the 3B), judge fluency bias (\S6.1). \\
J & Presentation: hardware framing, style, spelling & R1.3, R1.5, R3.5, R2 minors & ``Local server-side inference'' framing; Springer figure-reference style; run-to-run variance appendix; full proofread. \\
K & Related work additions & R1.6 & Two suggested papers added with substantive discussion; two respectfully declined as out of scope (rationale under R1.6). \\
\bottomrule
\end{tabular}
\end{center}

\section{Summary of Additional Experiments}
\label{sec:newexp}

\begin{enumerate}[leftmargin=1.4em]
  \item \textbf{Commensurable faithfulness re-judging} (R2.2): all 2{,}613 core outputs re-scored with source code as the single judge reference, as the mean of $k{=}3$ independent draws.
  \item \textbf{Blind two-annotator human study} (R2.6, R1.4, R2 minors): the same 51-item stratified sample re-annotated under a strictly blind protocol (no metric values visible, per-annotator shuffling); inter-rater $r=0.62$, ICC(2,1)$=0.63$.
  \item \textbf{Second LLM judge} (R1.4): Qwen3~8B as a cross-family judge over all outputs ($k{=}3$).
  \item \textbf{Three-arm exemplar ablation} (R2.8): original static few-shot vs.\ persona-corrected static vs.\ dynamically retrieved structurally matched exemplars, all three models.
  \item \textbf{Retrieval-corpus ablation} (R3.1): Simple RAG re-run with a library-specific corpus (Keras/scikit-learn developer documentation; leakage-controlled), all three models.
  \item \textbf{FP16 vs.\ Q4 quantization ablation} (R2.5): NoRAG family at FP16 on Qwen3~8B.
  \item \textbf{Corrected-gate rerun} (R1.2) and \textbf{repeated-run variance} (R2 minor): Iterative Critique family under a robust verdict parser; three end-to-end repeats of representative strategies.
\end{enumerate}

%====================================================================
\section{Reviewer 1}
%====================================================================

\subsection*{R1.1 --- Qwen2.5 in the abstract vs.\ Qwen3 in the body; exact model tags}
\rcomment{Why does the abstract mention Qwen2.5, while the main text, method section, and results use Qwen3 8B? \dots specify the actual Ollama model tag used.}
\resp Corrected. ``Qwen2.5'' was a remnant of a pilot experiment wave; the reported experiments used Qwen3. The revised \S4.2 states the exact served tags: \texttt{llama3.2:latest} (Llama~3.2~3B), \texttt{qwen3:8b}, \texttt{phi4:14b}, with \texttt{deepseek-coder:6.7b} as helper/judge and \texttt{qwen3:8b-fp16} in the quantization ablation. Our tag audit also uncovered a more serious naming error, reported under R2.1.

\subsection*{R1.2 --- Fragility of the critique-gate substring check}
\rcomment{Iterative Critique RAG determines whether retrieval should be triggered solely by checking whether the output contains the string ``EXCELLENT.'' Is this mechanism overly fragile?}
\resp The concern is valid, and our code audit found the situation was worse than described: the manuscript misdocumented the gate. As run, the check was \texttt{if "GOOD" not in response.upper()} (case-insensitivity was therefore already handled, but label-list echoes and negated phrasings could indeed defeat it). We have (a) corrected the description (revised \S3.3, Appendix~H); (b) implemented a robust verdict parser (assessment-line extraction with label-list-echo and negation handling, unit-tested); and (c) re-run the Iterative Critique family under the corrected gate on all three models: on Phi-4 (all four modes) and Qwen3 (Base), quality metrics change by less than $\pm0.01$ BERTScore and API calls by less than $\pm0.2$ per sample; on Llama~3.2~3B the corrected gate modestly improves quality ($+0.036$ BERTScore) while the condition remains below the zero-shot baseline. The heuristic therefore did not materially affect any reported conclusion (revised \S5.7).

\subsection*{R1.3 --- A100 hardware vs.\ on-device claims}
\rcomment{\dots the experimental hardware is an A100 GPU, which does not represent real edge devices\dots}
\resp Accepted. All ``on-device'' claims are reframed as \emph{local server-side inference}---the models are deployable on commodity on-premises GPUs, but no end-user-device measurements were made---and the hardware appendix now describes the cloud A100 environment precisely (revised \S4.2, Appendix~F).

\subsection*{R1.4 --- DeepSeek-Coder as both critic and judge}
\rcomment{Could this setup introduce evaluation bias? If so, do the authors provide any mitigation strategy?}
\resp The reviewer's concern proved more consequential than either we or the reviewer anticipated. We added a second judge from a different family (Qwen3~8B, $k{=}3$ draws, all 2{,}613 outputs) and validated both judges against a strictly blind human study. Both judges \emph{fail} validation (DeepSeek: $r=-0.21$; Qwen3: $r=0.05$; excluding Qwen3's own generations: $r=-0.09$, i.e., no self-preference), while agreeing almost perfectly with each other on strategy rankings ($\rho=0.979$ over all 39 cells): they share a systematic bias toward retrieval-grounded fluency ($+0.2$--$0.3$ over-rating of retrieval/critique outputs relative to blind humans). The revision therefore removes LLM judges from all primary claims, rests conclusions on human-validated reference-based metrics (pydocstyle $r=0.56$, rescaled BERTScore $r=0.52$ against blind humans), and reports the judge findings as a contribution (revised \S5.1, \S7.1, Appendix~B). The structural concern the reviewer raised---critic/judge role overlap---is thereby resolved: no LLM judge underpins any primary result. From the revised \S5.1:
\excerpt{The judges share a systematic criterion that is not human-aligned: both over-rate retrieval- and critique-based outputs by $+0.2$--$0.3$ relative to blind humans\ldots rewarding fluent retrieval-grounded paraphrase that human raters penalize for loss of code-specific detail.}

\subsection*{R1.5 --- ``Fig.'' vs.\ ``Figure'' consistency}
\resp Corrected throughout to Springer style; a full cross-reference audit was performed.

\subsection*{R1.6 --- Suggested related work}
\resp We added and discuss the two suggestions most relevant to our study: calibration of language models on code summarization (Proc.\ ACM Softw.\ Eng., 2025), which directly supports our judge-reliability findings (\S2.5), and rethinking-based code summarization with a chain of comments (COLING 2025) in the reasoning-strategies discussion (\S2.1). We respectfully did not add the mixed-precision quantization (NeurIPS 2026) and PSO-driven medical information extraction (Swarm Evol.\ Comput., 2025) papers, whose topics (quantization algorithm design; medical IE) fall outside the scope of our related work; our quantization ablation cites the empirical motivation directly.

%====================================================================
\section{Reviewer 2}
%====================================================================

\subsection*{R2.1 --- Model identities}
\rcomment{Meta's Llama 3.2 release contains no 8B model\dots misidentifying one of the two compared models is disqualifying as written.}
\resp The reviewer is correct, and we apologise for the error. Auditing our execution logs established that the served model was \texttt{llama3.2:latest}, i.e., \textbf{Llama~3.2~3B}. All references are corrected; the study's true capacity range is \textbf{3B--14B}. The capacity finding was re-verified for the models actually used and \emph{strengthens} under the correction: under the revised mixed-effects analysis, Phi-4~14B and Qwen3~8B outperform Llama~3.2~3B on commensurable faithfulness by $+0.047$ and $+0.064$ respectively (both $p<0.001$) and on BERTScore by $+0.024$ and $+0.015$ (both $p<0.001$). The abstract's ``Qwen2.5'' is likewise corrected (see R1.1). We have additionally corrected the reported benchmark scale (see our note to the editor): the executed benchmark is 67 classes / 2{,}613 core generations, and all tables and statistics in the revision reflect the executed data.

\subsection*{R2.2 --- Faithfulness not commensurable across families}
\rcomment{Faithfulness should be judged against a single reference (the source code) for all strategies\dots}
\resp Fully accepted and implemented. Every output was re-judged against source code as the single reference (identical rubric; $k{=}3$-draw means; both judges). Under the commensurable metric and the corrected unit of analysis, no architecture family shows a significant faithfulness effect (Simple RAG $-0.011$, Iterative Critique $+0.006$, Few-Shot $-0.017$; all ns), GoT-style prompting is the only significant reasoning effect ($+0.032$, $p=0.003$), and model effects dominate. The context-referenced score is retained only as a supplementary grounding indicator (Appendix~C). This correction also resolved the inconsistency the reviewer identified in R2.6.

\subsection*{R2.3 --- Unit of analysis}
\rcomment{A mixed-effects model with class as a random effect\dots is the appropriate framework.}
\resp Fully accepted and implemented exactly as proposed. All inferential claims now derive from linear mixed-effects models over the $N{=}2{,}613$ per-class observations (class random intercept; family, mode, and model fixed effects; REML). The class ICC (up to 0.28) confirms the dependence structure the reviewer identified. Final estimates (revised Table~5): BERTScore --- Few-Shot $-0.080^{***}$, Iterative Critique $-0.034^{***}$, Simple RAG $-0.006$ ($p{=}.052$), GoT $-0.009^{*}$, Qwen3 $+0.015^{***}$, Phi-4 $+0.024^{***}$; faithfulness --- GoT $+0.032^{**}$, Qwen3 $+0.064^{***}$, Phi-4 $+0.047^{***}$, all family effects ns. Several aggregate-level contrasts from the submitted version did not survive and the text has been revised accordingly; the prior Kruskal--Wallis/Mann--Whitney apparatus is removed. From the revised \S3.7:
\excerpt{Because all 67 benchmark classes are evaluated under every model--strategy condition, per-condition scores are not independent across conditions. All inferential claims therefore derive from linear mixed-effects models fitted on the per-class observations ($N{=}2{,}613$ per metric), with the benchmark class as a random intercept and architecture family, reasoning mode, and model as fixed effects\ldots The class random intercept is substantial (ICC $=0.28$ for BERTScore), confirming that class identity induces the dependence this design accounts for.}

\subsection*{R2.4 --- Headlining a metric the paper discredits}
\resp Accepted. The revised paper leads with a metric-validation section (\S5.1) and rests all claims on human-validated metrics. BERTScore was additionally recomputed with \texttt{roberta-large} and baseline rescaling (R2 minor): strategy orderings are backbone-robust, and the rescaled absolute values (near zero) confirm the reviewer's point that the original absolute values were difficult to interpret. The GoT/BERTScore result is no longer headlined; it appears with its honest per-class effect size ($-0.009$).

\subsection*{R2.5 --- Quantization confound}
\rcomment{\dots the conclusion cannot be separated from ``4-bit quantization suppressed the reasoning benefit.''}
\resp We ran the requested ablation: the NoRAG family at FP16 vs.\ Q4\_K\_M on Qwen3~8B. Deltas are uniformly an order of magnitude below the effects of interest (Plain $-0.004$, CoT $+0.005$, GoT $+0.007$, ToT $+0.004$ BERTScore), while FP16 substantially increases latency. The reasoning-mode conclusions are therefore not quantization artifacts; the confound is discussed and dismissed with data (revised \S5.4).

\subsection*{R2.6 --- Human validation contradicts the main results}
\label{resp:r26}
\rcomment{The claim that ``human evaluation confirms our automated rankings'' is not supported by the numbers beneath it and must be reconciled or withdrawn.}
\resp The reviewer identified a real inconsistency whose root cause proved deeper than the reference mismatch of R2.2. Investigating it, we found (a) per-sample LLM-judge scores at sampling temperature have near-zero test--retest reliability (identical inputs score 0.5--1.0 across draws), which is mathematically incompatible with the reported $r=0.925$ human agreement; and (b) the original annotation sheet displayed the judge's scores beside the entry field. A strictly blind re-annotation of the same 51 items by two annotators (metric values hidden, order shuffled per annotator, key sealed) correlates only $r=0.16$ with the original ratings and $r=0.13$ with the original judge scores, indicating the original ratings were anchored. We have withdrawn the $r=0.925$ claim and the ``human evaluation confirms our automated rankings'' statement, replaced Section~3.6/Table~6/Fig.~9 with the blind study (inter-rater ICC $=0.63$), and rebuilt the paper's metric hierarchy on it (revised \S3.6, \S5.1, \S7.1). The blind study also reverses the specific ordering the reviewer questioned: blind humans rank Plain+GoT highest (0.71) and Simple RAG fourth (0.57). We report the full episode in Threats to Validity as a caution for LLM-as-judge methodology. From the revised \S7.1:
\excerpt{\ldots a strictly blind re-annotation of the same 51 items by two annotators correlates only weakly with those original ratings ($r=0.16$) and with the judge ($r=0.13$), indicating the original ratings were anchored to the visible scores. We therefore withdrew the original validation and rebuilt all human-referenced claims on the blind study. We report this experience as a caution for the growing LLM-as-judge literature: agreement estimates obtained under non-blind protocols can be arbitrarily inflated, and per-sample judge scores at sampling temperature are draw-noise unless averaged.}

\subsection*{R2.7 --- Canonical ToT/GoT vs.\ implemented variants}
\resp Accepted. The modes are now described as \emph{single-pass, prompt-level adaptations} (``Reasoning-Mode Prompt Variants,'' revised \S3.4) with an explicit scoping statement; labels are retained to indicate lineage, and all conclusions are scoped to the implemented variants.

\subsection*{R2.8 --- Few-Shot baseline too weak}
\rcomment{Dynamic or structurally matched exemplar retrieval is the standard baseline and should be included before ``Few-Shot is harmful'' is stated as a general finding.}
\resp Accepted and implemented, with an additional disclosure: code review found the static few-shot prompt contained a conflicting evaluator persona, which we disclose in \S3.3. The revision reports a three-arm ablation on all three models: original static, persona-corrected static, and dynamically retrieved structurally matched exemplars (leave-one-out nearest neighbors with human reference docstrings). The result overturns our own submitted finding: \textbf{dynamic exemplars outperform the zero-shot baseline on every model} ($+0.041$ to $+0.087$ BERTScore; ROUGE-1 roughly doubles), producing the strongest condition in the study, while the persona correction alone interacts with capacity (rescuing Phi-4, further degrading the 3B model). ``Few-Shot is harmful'' is withdrawn and replaced by: static, structurally mismatched exemplars are harmful; matched exemplar retrieval is the study's most effective intervention (revised \S5.5, \S6.2). From the revised \S6.2:
\excerpt{This reframes the practical question for retrieval-augmented documentation tooling from ``should we add RAG?'' to ``what should the index contain?''---and our evidence says: exemplar pairs, not style guides.}

\subsection*{R2.9 --- ``Scaling'' framing}
\resp Accepted; we adopt the reviewer's own wording. The title is now ``\emph{Capacity} vs.\ Architecture\dots'', and the text attributes model effects to capacity \emph{and} training quality, noting explicitly that a three-vendor design cannot separate the two (revised \S5.2, \S7.2).

\subsection*{R2 minor points}
\resp All addressed: BERTScore backbone and rescaling (see R2.4); coverage metrics now descriptive-only with an explicit statement (\S3.5, \S5); the $\rho=-0.286$ divergence between faithfulness measures is reinterpreted mechanistically---the prior claim that it ``validates'' the judge is removed, and token-overlap is repositioned as a rewriting-depth indicator (Appendix~C); the forward-dated non-archival motivational source is removed from the introduction; the ambiguous ``N'' column is gone---the underlying table was superseded by the mixed-effects presentation; repeated runs quantify run-to-run variance (strategy-mean BERTScore SDs of $0.0015$--$0.0017$, one to two orders below the effects of interest; Appendix~E), which also addresses the stability of cost figures following the disclosed ToT API-count correction.

%====================================================================
\section{Reviewer 3}
%====================================================================

\subsection*{R3.1 --- Retrieval corpus fairness}
\rcomment{\dots the retrieval corpus mainly consists of general knowledge\dots whether the comparison is fair remains questionable.}
\resp A fair concern, tested directly: we re-ran Simple RAG with a \emph{library-specific} corpus built from the benchmark libraries' developer guides and convention documentation (leakage-controlled: per-class API reference pages excluded), on all three models. Corpus relevance matters directionally---the library corpus beats the generic corpus on every model and repairs the parameter-coverage damage that generic retrieval caused on the smallest model (0.34$\rightarrow$0.51)---but never exceeds the zero-shot baseline by more than $+0.012$ BERTScore. The retrieval conclusion is therefore retained in properly scoped form: documentation-context retrieval is corpus-sensitive but marginal, whereas exemplar retrieval (R2.8) is highly effective---``what you retrieve matters more than whether you retrieve'' (revised \S5.3, \S6.2).

\subsection*{R3.2 --- Reasoning Tax lacks mechanism}
\resp The revision adds a mechanism section (\S6.1) with three quantified failure modes: \emph{hallucination by inference} (speculative parametric associations, illustrated with the Keras example); \emph{capacity-dependent format-compliance failures} (all five runaway generations of $2{,}000$--$290{,}000$ words occurred in the 3B model's reasoning modes; 36\% degenerate truncation under static few-shot on Qwen3); and \emph{judge fluency bias} (why automated scores diverge from human judgment on refined outputs). The Reasoning Tax framing is correspondingly refined: GoT-style enumeration carries a real, human-corroborated faithfulness benefit at $3.1\times$ cost, while ToT-style branching remains a pure tax.

\subsection*{R3.3 --- Novelty vs.\ existing benchmarks}
\resp The contribution statement is sharpened (\S1, \S2.6): cross-model factorial design at class level, five targeted ablations, cost-aware Pareto analysis, and---unique in this space to our knowledge---a strictly blind, two-annotator validation of every automated metric, which materially changed the paper's conclusions and which we release for reuse.

\subsection*{R3.4 --- Scale vs.\ training-data confound}
\resp Accepted; resolved jointly with R2.9. Claims are rescoped to capacity/training quality, and the limitation is stated explicitly with within-family comparisons identified as future work (\S7.2).

\subsection*{R3.5 --- Spelling and grammar}
\resp A full proofreading pass was performed on the revised manuscript.

\bigskip
\noindent Thank you again for the time and care invested in this manuscript. The review process materially improved this work---several of the revision's strongest results exist because the reviewers insisted on them. We are happy to provide additional clarifications, supplementary analyses, or annotated artifacts that would assist the second-round review.

\end{document}
